\documentclass[11pt]{article}

\usepackage[margin=1in]{geometry}
\usepackage{amsmath}
\usepackage{amssymb}
\usepackage{amsthm}
\usepackage{enumitem}

\theoremstyle{definition}
\newtheorem{definition}{Definition}

\theoremstyle{plain}
\newtheorem{lemma}{Lemma}
\newtheorem{proposition}{Proposition}

\title{Numbered Increasing $(1,2)$-Trees and the Euler Zigzag Numbers}
\author{Directed Reading Program, Northwestern University\\[2pt]
\small In collaboration with Noah Wisdom (graduate mentor)}
\date{Spring 2024}

\begin{document}
\maketitle

\section*{Some background}

The \emph{Euler numbers}, or \emph{Euler zigzag numbers}, are the sequence
\[
A_n = 1,\ 1,\ 1,\ 2,\ 5,\ 16,\ 61,\ 272,\ \dots
\]
They satisfy the recurrence
\[
2A_{n+1} \;=\; \sum_{k=0}^{n} \binom{n}{k} A_k\, A_{n-k}.
\]
This comes up a lot in combinatorics.

\bigskip

The objects we will count are the following.

\begin{definition}[Numbered increasing $(1,2)$-trees]
A \emph{numbered increasing $(1,2)$-tree} is an increasing rooted tree such that
every non-endpoint vertex has one or two children. Two trees on the same number
of vertices are considered the same (isomorphic) if one can be reoriented so as
to recover the other. The trees are numbered (labeled $1, 2, \dots, n$) such that
the labels strictly increase along every path from the root to a leaf.
\end{definition}

The notes illustrate this with two small examples: a pair of trees that are
considered the same, since one can be rearranged to match the other, and a pair
of trees that are distinct, since no such rearrangement exists.

\bigskip

We write
\[
f(n) = \text{the number of distinct numbered increasing $(1,2)$-trees on $n$ vertices.}
\]

\begin{lemma}\label{lem:egf}
\[
\sum_{n \ge 0} E_n\, \frac{x^n}{n!} \;=\; \sec x + \tan x .
\]
\end{lemma}

\section*{Main result}

\begin{proposition}
For each $n$, $\; f(n) = E_n.$
\end{proposition}

Before the proof, we check the first few values directly:
\[
f(1) = 1, \qquad f(2) = 1, \qquad f(3) = 2, \qquad f(4) = 5,
\]
which agree with $E_1, E_2, E_3, E_4$.

\begin{proof}
Let
\[
y \;=\; \sum_{n \ge 1} f(n)\, \frac{x^n}{n!}
   \;=\; x + \frac{x^2}{2!} + 2\,\frac{x^3}{3!} + \cdots .
\]
Taking the derivative with respect to $x$ and reindexing,
\[
y' \;=\; \sum_{n \ge 1} f(n)\, \frac{n\,x^{\,n-1}}{n!}
   \;=\; \sum_{n \ge 0} f(n+1)\, \frac{x^n}{n!}.
\]
So $y'$ is the exponential generating function counting numbered increasing
$(1,2)$-trees on $n+1$ vertices. We split such a tree into cases according to
the structure at the root:
\begin{enumerate}[label=(\alph*)]
  \item a single vertex $(n = 0)$;
  \item the tree has one subtree of the root, which is an increasing
        $(1,2)$-tree with $n$ vertices;
  \item the tree has two subtrees of the root, each of which is an increasing
        $(1,2)$-tree, with $n$ vertices in all.
\end{enumerate}
From these cases we obtain
\[
y' \;=\; 1 + y + \tfrac{1}{2} y^2 .
\]
(This checks out if you compute it, but we will not carry out that computation
here.)

\medskip

This differential equation has the solution
\[
y \;=\; \sec x + \tan x - 1 .
\]
To see why, compute the derivative
\[
y' \;=\; \sec x \tan x + \sec^2 x ,
\]
and substitute $y$ into the right-hand side of the equation:
\begin{align*}
1 + y + \tfrac{1}{2} y^2
 &= 1 + \bigl(\sec x + \tan x - 1\bigr)
      + \tfrac{1}{2}\bigl(\sec x + \tan x - 1\bigr)^2 \\[2pt]
 &= \bigl(\sec x + \tan x\bigr)
      + \tfrac{1}{2}\bigl(\sec^2 x + \tan^2 x + 1
        + 2\sec x \tan x - 2\sec x - 2\tan x\bigr) \\[2pt]
 &= \sec x \tan x + \tfrac{1}{2}\bigl(\sec^2 x + \tan^2 x + 1\bigr) \\[2pt]
 &= \sec x \tan x + \tfrac{1}{2}\bigl(\sec^2 x + \sec^2 x\bigr) \\[2pt]
 &= \sec x \tan x + \sec^2 x \\[2pt]
 &= y',
\end{align*}
where we used $\tan^2 x + 1 = \sec^2 x$. This is exactly what we wanted.

\medskip

Therefore
\[
y \;=\; \sum_{n \ge 1} E_n\, \frac{x^n}{n!},
\]
which is the generating function of Lemma~\ref{lem:egf} with its constant term
removed. Matching coefficients gives
\[
f(n) = E_n \qquad (\text{starting from } E_1). \qedhere
\]
\end{proof}

\end{document}
